\section{本章小结}
\label{sec:ns-summary}

本章从一个具体的麻烦开始：Lasso的解恰好落在目标不可微的地方。为此先用支撑直线的斜率定义了次梯度和次微分，凸目标的最优解由$\bm{0}\in\partial f$充要地刻画；和与复合的规则需要相对内部这样的正则条件。非凸函数上要改用局部定义的广义导数，那时的条件只是必要的。

此后各节回答“用什么信息推进迭代”，答案不同，要求和保证也不同。

次梯度法只用支撑信息，每步一个次梯度，单步不一定下降，但距离的递推使平均点或最好点的目标差为$O(1/\sqrt{T})$，这在黑箱模型下不能本质改进。

邻近方法把难求导的结构留在子问题里。复合目标的邻近梯度法（ISTA）在凸情形达到$O(1/T)$，FISTA用外推达到$O(1/T^2)$，后者不要求逐轮下降，也不能直接用于非凸目标；非凸情形用映射残差评价，残差为零只说明是复合驻点。邻近方法的前提是邻近映射可算，映射可定义、子问题可计算和原问题收敛是三件要分别核对的事。

光滑化和代理模型改变每一步要解的问题。Huber函数是绝对值的Moreau包络，Nesterov平滑利用最大值结构在内层加强凸项，它们用偏差$\mu D$换取$O(1/\mu)$的光滑常数，原目标的误差由偏差和求解误差两部分决定；平滑改变目标，加速改变迭代规则，配合使用可得$O(1/\varepsilon)$的总步数。MM用上界代理保证原目标不增，束方法用支撑线拼成的下界模型预测下降，并用接受规则和预测下降值来评价真实进展。

\begin{BookInlineTable}
  \centering
  \small
  \begin{tabular}{@{}lll@{}}
    \toprule
    方法 & 需要的结构 & 凸情形的保证\\
    \midrule
    次梯度法 & 次梯度有界 & 目标差$O(1/\sqrt{T})$\\
    邻近点法 & $\operatorname{prox}_{\alpha f}$可算 & 目标差$O(1/\sum\alpha^{(t)})$\\
    ISTA & $\nabla f$，$\operatorname{prox}_{\alpha r}$ & 目标差$O(1/T)$，单调\\
    FISTA & 同上 & 目标差$O(1/T^2)$，不单调\\
    平滑加速 & 最大值结构，内层可解 & $O(1/\varepsilon)$步，带偏差$\mu D$\\
    MM & 上界代理 & 目标值不增\\
    束方法 & 次梯度，二次规划求解 & 可计算的误差证书\\
    \bottomrule
  \end{tabular}
  \caption{本章方法的对照。$T$为迭代步数，$\varepsilon$为原目标精度。}
  \label{tab:ns-summary}
\end{BookInlineTable}

表~\ref{tab:ns-summary}把这些方法放在一起。本章全部在无约束或简单约束下讨论，一般约束的乘子、增广Lagrange函数与ADMM由第~\ref{chap:penalty-multiplier-methods}章结合约束结构建立；按变量分块使用邻近映射的方法见第~\ref{chap:block-optimization}章。
